\scanpage{75}
% slanted
\reversemarginpar\margin{\shortstack[l]{\unsure{joindre}\\SGA 7 \(\overline{\underline{\mathrm{VI}}}\) ?\\App. \(\overline{\underline{\mathrm{II}}}\)}}\normalmarginpar
\noindent\hfill Bures sur Yvette, le 21 octobre 1968\\

\noindent Cher Grothendieck,\\

En fait, il n'est pas nécessaire de faire de savants recollages de complexes cotangents pour prouver que

\noindent \uline{proposition} : \uline{Soit} \(X\) \uline{un espace analytique complexe compact}. \uline{Pour que} \(X\) \uline{soit un schéma}, \uline{il faut et suffit que} \(X_{\mathrm{red}}\) \uline{le soit}

On se ramène à \deleted{supposer} montrer que si \(X_0\) est défini par un Idéal de carré nul, soit \(\mathscr{N}\), de \(X\), alors
\[
 (X_0\text{ est un schéma})\ \Longrightarrow\ (X\text{ est un schéma})
\]

Soit \(E^{\mathrm{an}}(X_0,\mathscr{N})\) et \(E^{\mathrm{alg}}(X_0,\mathscr{N})\) les groupes analytiques et algébriques d'extensions infinitésimales de \(X_0\) par \(\mathscr{N}\). Soit enfin \(\unsure{K|\mathbb{K}}\) le complexe cotangent de \(X_0\), \uline{donné localement, dans la catégorie dérivée}. il s'obtient en plongeant \uline{localement} (Zariski, ou étale) \(X_0\) dans un \unsure{schéma} lisse \ldots\ et permet de définir
\begingroup
\makeatletter\tagsleft@true\makeatother
\[
 \mathscr{H}om(K,\mathscr{N})\qquad\text{et}\qquad\mathscr{E}xt^{1}(K,\mathscr{N})\ .\tag*{(1)}
\]
\endgroup

Soit le diagramme :

% Detached pair beside the upper diagram row; punctuation value unresolved.
\margin{\vspace*{1.4em}\figure{\unsure{%
\begin{tikzpicture}[x=1mm,y=1mm]
\fill (-0.22,-0.07)
 .. controls (-0.27,0.06) and (-0.16,0.16) .. (-0.08,0.17)
 .. controls (-0.04,0.30) and (0.06,0.25) .. (0.11,0.14)
 .. controls (0.25,0.10) and (0.27,-0.03) .. (0.15,-0.11)
 .. controls (0.04,-0.15) and (-0.15,-0.16) .. (-0.22,-0.07) -- cycle;
\fill (0.37,0.72)
 .. controls (0.34,0.85) and (0.44,0.95) .. (0.50,0.98)
 .. controls (0.54,1.10) and (0.65,1.11) .. (0.71,1.00)
 .. controls (0.83,0.94) and (0.77,0.85) .. (0.68,0.81)
 .. controls (0.73,0.69) and (0.64,0.62) .. (0.53,0.65)
 .. controls (0.44,0.63) and (0.38,0.65) .. (0.37,0.72) -- cycle;
\end{tikzpicture}%
}}}

\diagram{\footnotesize
\begin{tikzcd}[ampersand replacement=\&,column sep=0.85em,row sep=1.6em]
 (2) \& 0 \arrow[r] \&
 H^{1^{\mathrm{alg}}}(X_0,\mathscr{H}om(K,\mathscr{N})) \arrow[r] \arrow[d] \&
 E^{\mathrm{alg}}(X_0,\mathscr{N}) \arrow[r,"\alpha"] \arrow[d] \&
 H^{0}(X_0,\mathscr{E}xt^{1}(K,\mathscr{N})) \arrow[d] \\
 \& 0 \arrow[r] \&
 H^{1^{\mathrm{an}}}(X_0,\mathscr{H}om(K,\mathscr{N})) \arrow[r] \&
 E^{\mathrm{an}}(X_0,\mathscr{N}) \arrow[r,"\beta"] \&
 H^{0^{\mathrm{an}}}(X_0,\mathscr{E}xt^{1}(K,\mathscr{N}))
\end{tikzcd}}


Vu les significations (locales) de (1), les lignes sont exactes ; par GAGA, les flèches verticales extrêmes sont des isomorphismes. Reste à \deleted{comparer les images de \mbox{\(\alpha\)} et \mbox{\(\beta\)}} montrer que \(\operatorname{Im}(\beta)\subset\operatorname{Im}(\alpha)\).
% archive sheet number: 74
\clearpage
